% Auto-generated from meeting_2567-02-16_mko2-finalization.md (กบ.วช. format v2)
\documentclass{meetingminutes}
\meetingcommittee{คณะกรรมการบริหารหลักสูตร วิศวกรรมคอมพิวเตอร์และปัญญาประดิษฐ์ (บัณฑิตศึกษา)}
\meetingnumber{๒}{๒๕๖๗}
\meetingdate{วันที่ ๑๖ กุมภาพันธ์ พ.ศ. ๒๕๖๗}
\meetingplace{ห้องประชุมคณะเทคโนโลยีอุตสาหกรรม}
\meetingstart{๑๓.๐๐ น.}
\meetingend{๑๕.๑๐ น.}

\begin{document}
\maketitle

\psection{ผู้มาประชุม}
\begin{participants}
  \pmember{1}{รองศาสตราจารย์ ดร.อิสระ อินจันทร์}{ที่ปรึกษาหลักสูตร}{ประธาน}
  \pmember{2}{รองศาสตราจารย์ ดร.กันต์ อินทุวงศ์}{ที่ปรึกษาหลักสูตร}{รองประธาน}
  \pmember{3}{ผู้ช่วยศาสตราจารย์ ดร.พิศิษฐ์ นาคใจ}{อาจารย์ประจำหลักสูตร}{กรรมการ}
  \pmember{4}{ผู้ช่วยศาสตราจารย์ ดร.กาญจนา ดาวเด่น}{อาจารย์ประจำหลักสูตร (เลขานุการ)}{กรรมการ/เลขานุการ}
  \pmember{5}{ผู้ช่วยศาสตราจารย์ ดร.พัชรี มณีรัตน์}{อาจารย์ประจำหลักสูตร}{กรรมการ}
  \pmember{6}{ผู้ช่วยศาสตราจารย์ ดร.พิทักษ์ คล้ายชม}{รองคณบดีฝ่ายวิชาการ คณะเทคโนโลยีอุตสาหกรรม}{กรรมการ}
  \pmember{7}{ผู้ช่วยศาสตราจารย์ ดร.โสกณ วิริยะรัตนกุล}{อาจารย์ประจำหลักสูตร}{กรรมการ}
  \pmember{8}{ผู้ช่วยศาสตราจารย์ ดร.ธัชชัย อยู่ยิ่ง}{อาจารย์ประจำหลักสูตร}{กรรมการ}
  \pmember{9}{ผู้ช่วยศาสตราจารย์ ดร.อภิศักดิ์ พรหมฝ่าย}{อาจารย์ประจำหลักสูตร}{กรรมการ}
\end{participants}

\psection{ผู้ไม่มาประชุม}
\noindent ไม่มี\par

\startmeeting
\openingchair{รองศาสตราจารย์ ดร.อิสระ อินจันทร์}{กล่าวเปิดการประชุมและดำเนินการประชุมตามระเบียบวาระ ดังนี้}

\agenda{๑}{รับรองรายงานการประชุมครั้งที่ 1/2567}
ที่ประชุมรับรองรายงานการประชุมครั้งที่ 1/2567 โดยไม่มีการแก้ไข


\agenda{๒}{ตรวจสอบ มคอ.2 ฉบับสมบูรณ์}
ผศ.ดร.พิศิษฐ์ นำเสนอ มคอ.2 ฉบับสมบูรณ์ที่จัดทำตามกรอบ กมอ. 2565 ที่ประชุมตรวจสอบความสมบูรณ์และความถูกต้องของเนื้อหาในแต่ละส่วน:

ส่วนที่ 1: ข้อมูลทั่วไป — ครบถ้วน ส่วนที่ 2: อาจารย์ผู้รับผิดชอบหลักสูตร — ระบุ 7 คน ครบถ้วน ส่วนที่ 3: สรุปผลการเรียนรู้ที่คาดหวัง (PLOs) — 5 PLOs พร้อม Bloom's Level ส่วนที่ 4: ข้อมูลเฉพาะของหลักสูตร — โครงสร้างรายวิชาและแผนการศึกษา ส่วนที่ 5: หลักเกณฑ์การประเมิน — เกณฑ์การวัดผลและการสำเร็จการศึกษา ส่วนที่ 6: การพัฒนาคณาจารย์ — แผนพัฒนาอาจารย์ประจำหลักสูตร ส่วนที่ 7: ตัวบ่งชี้คุณภาพและการประเมิน — KPIs และแผนติดตาม

ที่ประชุมพิจารณาแล้วมีข้อเสนอแนะให้ปรับปรุงถ้อยคำในส่วนที่ 3 (PLO3) ให้ระบุความแตกต่างระหว่างแผน 1 และแผน 2 ชัดเจนขึ้น

\resolution{เห็นชอบ มคอ.2 ฉบับสมบูรณ์ (หลังปรับปรุง) และอนุมัติให้เสนอต่อคณะกรรมการบัณฑิตศึกษาประจำมหาวิทยาลัย}

\agenda{๓}{กำหนดการสำรวจผู้มีส่วนได้ส่วนเสียอย่างเป็นทางการ (Formal Stakeholder Survey)}
ที่ประชุมเห็นชอบให้ดำเนินการสำรวจผู้มีส่วนได้ส่วนเสียอย่างเป็นทางการ โดยใช้ Google Form เพื่อเก็บข้อมูลเชิงปริมาณ กำหนดดำเนินการในช่วงเดือนมิถุนายน 2567

\resolution{มอบหมายให้ ผศ.ดร.กาญจนา ดาวเด่น ออกแบบแบบสอบถามและดำเนินการเก็บข้อมูล}

\adjournmeeting

\begin{sigblock}
\typist{ผู้ช่วยศาสตราจารย์ ดร.กาญจนา ดาวเด่น}
\reviewer{รองศาสตราจารย์ ดร.อิสระ อินจันทร์}{กรรมการ}
\chairsig{รองศาสตราจารย์ ดร.อิสระ อินจันทร์}{ประธานการประชุม}
\end{sigblock}

\end{document}
